\section{Contexte}

La HEIG-VD dispose d'un cluster \gls{kubernetes} équipé de \glspl{gpu}.
Dans le cadre de projets d'analyse d'images géospatiales, ce cluster doit supporter
un pipeline d'annotation automatique capable de traiter de grands volumes d'images haute résolution.

Les images sont stockées sur un serveur \gls{minio} (NAS Synology, protocole \gls{s3})
et les annotations sont gérées via \gls{labelstudio}.
Le modèle \gls{sam3} assure la segmentation automatique.

\section{Problématique}

Les principales contraintes sont :

\begin{itemize}
    \item les images haute résolution dépassent la capacité mémoire d'un seul nœud ;
    \item \gls{sam3} requiert un \gls{gpu} pour fonctionner à une vitesse acceptable ;
    \item les résultats de segmentation doivent être persistés et accessibles pour traitement ultérieur ;
    \item le pipeline doit être scalable sur le cluster existant ;
    \item le pipeline doit exposer des métriques pour surveiller l'utilisation du matériel et des modèles.
\end{itemize}

\section{Périmètre du travail}

Ce document définit le périmètre, les objectifs, les exigences et l'architecture du pipeline.
Il constitue la référence entre l'étudiant et le superviseur.
